\section{The inductive framework}\label{sec:induction}

The induction has two outputs in each dimension: smooth canonical
nonvanishing and good minimal models for real-boundary lc pairs.
This section proves the passage from the first output to the second,
assuming good models in smaller dimensions. All varieties here are
projective over \(\C\).

For an effective real boundary \(\Delta\), a good log minimal model of
\((X,\Delta)\) is a log minimal model on which the adjoint divisor is
semiample. For real divisors, semiampleness means real linear equivalence
to the pullback of an ample real divisor by a contraction. We allow the
usual definition of a log minimal model in which extracted prime
divisors are included in its boundary with coefficient one.
For a pseudo-effective divisor \(A\), we use
\(\kappa_\sigma(X,A)\) for Nakayama's numerical dimension. When \(A\)
is nef this agrees with the intersection-theoretic numerical dimension
\(\nu(X,A)\).

Fix an integer \(n\geq 1\).
\begin{assumption}[Lower-dimensional good models]\label{mmp:lower-models}
Every projective log canonical pair of dimension less than \(n\),
with real boundary and pseudo-effective adjoint divisor, has a good
log minimal model.
\end{assumption}

The base case is a point. At the next step, smooth nonvanishing in
dimension \(n\) will be proved from Assumption~\ref{mmp:lower-models}
in Theorem~\ref{thm:smooth-nonvanishing}. It will then become an input
to Proposition~\ref{prop:inductive-reduction}. The boundary-restriction
argument below and the special-termination argument in
Section~\ref{sec:exclusions} are available before this same-dimensional
nonvanishing step; they use only the stated lower-dimensional hypothesis.

For clarity, let \(\mathrm{GLM}_{\le d}\) denote good-model existence
for all projective complex lc pairs with effective real boundary and
pseudo-effective adjoint in dimensions at most \(d\). Let
\(\mathrm{NV}_n\) denote smooth canonical nonvanishing in dimension
\(n\), and let \(\mathrm{SA}_n\) denote semiampleness of nef rational
lc adjoints in that dimension. The proof proceeds in the following
order; the real-boundary conclusion in the third row supplies the next
dimension's hypothesis.
\begin{center}
\small
\begin{tabular}{@{}p{.19\textwidth}p{.37\textwidth}p{.35\textwidth}@{}}
\toprule
Stage & Input & Output and location \\
\midrule
Base & Dimension zero & \(\mathrm{GLM}_{\le0}\): a point \\
Smooth step & \(\mathrm{GLM}_{\le n-1}\) and
Theorem~\ref{asm:iitaka} & \(\mathrm{NV}_n\),
Theorem~\ref{thm:smooth-nonvanishing} \\
Good-model step & \(\mathrm{GLM}_{\le n-1}\) and \(\mathrm{NV}_n\)
& \(\mathrm{GLM}_{\le n}\),
Proposition~\ref{prop:inductive-reduction} \\
Nef conclusion & \(\mathrm{GLM}_{\le n}\) & \(\mathrm{SA}_n\),
Lemma~\ref{mmp:comparison} \\
\bottomrule
\end{tabular}
\end{center}
The signed-representative theorem used in both steps is proved in
Sections~\ref{sg:section}--\ref{sg:descent}; its whole-boundary
semiampleness hypothesis comes from the lower-dimensional row,
via Proposition~\ref{prop:boundary-restriction}.

\subsection{Comparison and boundary restrictions}

We begin with two forms of descent. The first compares the adjoint on
a pair and its minimal model as actual divisors on a common resolution.
The second joins the sections obtained by adjunction on the components
of a reduced dlt boundary.

\begin{lemma}[Comparison with a nef model]\label{mmp:comparison}
Let \((X,\Delta)\) be log canonical and let \((X',\Delta')\)
be a log minimal model. On a common resolution
\[
 X \xleftarrow{u} W \xrightarrow{v} X'
\]
there is an effective \(v\)-exceptional divisor \(F\) such that
\begin{equation}\label{mmp:comparison-equation}
 u^*(K_X+\Delta)=v^*(K_{X'}+\Delta')+F.
\end{equation}
If \(K_X+\Delta\) is nef, then \(F=0\). Consequently a nef
rational adjoint is semiample whenever it has a good log minimal
model. If \((X,\Delta)\) is klt, its log minimal model is klt and
extracts no divisors.
\end{lemma}

\begin{proof}
Put \(F=u^*(K_X+\Delta)-v^*(K_{X'}+\Delta')\).
Discrepancy improvement on divisors of \(X\) gives \(u_*F\geq 0\),
while \(-F\) is \(u\)-nef because \(K_{X'}+\Delta'\) is nef.
The negativity lemma therefore gives \(F\geq 0\).
At a prime of \(W\) that is not \(v\)-exceptional, the coefficient
of \(F\) is zero unless that prime is extracted on \(X'\).
In the latter case it equals the negative of its log discrepancy
over \((X,\Delta)\). It is therefore nonpositive, and hence zero.
Thus \(F\) is \(v\)-exceptional. If the source adjoint is nef,
then \(F\) is also \(v\)-nef, so the negativity lemma gives \(F=0\).

The equality of pullbacks transfers semiampleness of a rational
adjoint. Indeed a proper birational morphism onto a normal variety
has direct image of its structure sheaf equal to that structure
sheaf. Projection formula therefore identifies the global sections
of each Cartier multiple with the sections of its pullback.
If every such pulled-back section vanished over a point downstairs,
it could not generate upstairs over that point. Generation upstairs
thus implies generation downstairs.
Finally, for a klt source the coefficient at an extracted prime
would be strictly negative, which is impossible. The remaining
discrepancy inequalities show that the model is klt.
\end{proof}

\begin{proposition}[Semiampleness on the reduced boundary]
\label{prop:boundary-restriction}
Assume Assumption~\ref{mmp:lower-models}. Let \((V,D)\) be a
projective \(\Q\)-factorial dlt \(n\)-fold with \(D\) reduced and
\(M=K_V+D\) nef. Then \(M|_D\) is semiample as a rational line
bundle on the reduced scheme \(D\).
\end{proposition}

\begin{proof}
The ambient variety \(V\) is klt. The reduced floor of a
\(\Q\)-factorial dlt pair is \(S_2\), and it has ordinary double
crossings in codimension one; its irreducible components are normal.
For the \(S_2\) assertion one can work locally with the cyclic
class cover of the \(\Q\)-Cartier divisor \(D\). This cover is
quasi-\'etale and klt, hence Cohen--Macaulay. The eigensheaf
\(\OO_V(-D)\), being a direct summand of its finite direct image,
is Cohen--Macaulay. The divisor sequence then shows that \(D\)
is Cohen--Macaulay.

Divisorial adjunction on the normalization
\(\coprod D_i\to D\) gives lc pairs
\((D_i,\operatorname{Diff}_{D_i}(D-D_i))\).
The different includes the conductor with coefficient one.
Removing that conductor contribution defines an effective boundary
on \(D\); together with \(D\) it is an slc pair whose adjoint
line bundle is \(M|_D\). The divisible residue identifications
hold in codimension one, with even powers eliminating residue
signs at double crossings, and extend by \(S_2\).

Each normalized adjoint is nef and semiample by
Assumption~\ref{mmp:lower-models} and
Lemma~\ref{mmp:comparison}.
Semiampleness descends from the normalization by the slc gluing
theorem \cite[Theorem~1.5, equivalently Theorem~4.3]
{FujinoGongyo2014}. This gives the required statement on the
whole reduced scheme, not merely on its individual components.
\end{proof}

\subsection{A uniform trivial-perturbation argument}

We will also perturb a nef adjoint into a klt program. A single Cartier
index must survive every step: otherwise the smallness required of the
perturbation could change along the program. The following estimate and
line-bundle descent give that uniformity.

\begin{lemma}\label{mmp:trivial-perturbation}
Let \((Z,\Delta)\) be a projective \(\Q\)-factorial dlt rational
pair of dimension $n$, and suppose \(L=K_Z+\Delta\) is nef.
Set \(P=\Delta\) if the pair is klt and
\(P=\lfloor\Delta\rfloor\) otherwise.
Fix \(m>0\) such that \(mL\) is Cartier.
For every rational
\[
 0<\epsilon<\frac{1}{4nm+2},
\]
every step of an LMMP for \(L-\epsilon P\) is \(L\)-trivial.
On all its models, the transform of \(L\) is nef and its
multiple by the same integer \(m\) is Cartier.
\end{lemma}

\begin{proof}
Both \(\Delta-\epsilon P\) and \(\Delta-\tfrac12P\) are
effective klt boundaries. If \(R\) is a ray negative for
\(L-\epsilon P\), nefness of \(L\) shows that \(R\) is also
negative for \(L-\tfrac12P\).
Choose a rational curve \(C\) spanning this ray and satisfying
the length bound
\[
 -\bigl(L-\tfrac12P\bigr)\cdot C\leq 2n.
\]
Writing \(a=L\cdot C\geq 0\) and \(p=P\cdot C\), we have
\[
 p\leq 2a+4n,\qquad a<\epsilon p,
 \qquad
 (1-2\epsilon)a<4n\epsilon.
\]
Our choice of \(\epsilon\) gives \(a<1/m\).
Since \(ma\) is a nonnegative integer, \(a=0\).

For the driving klt contraction $c:Z\to B$, the line bundle
\(\OO_Z(mL)\) is numerically trivial over $B$.
The line-bundle clause of the contraction theorem gives its
descent to a line bundle $\mathcal A_B$ on $B$, with no change of \(m\)
\cite[Theorem~3.74(3)]{Fujino2017}.
One may also see the unchanged index directly from relative
basepoint freeness: two consecutive sufficiently large powers
descend, so their quotient descends.
The line bundle $\mathcal A_B$ is nef because its pullback
$\OO_Z(mL)$ is nef. For a flip $c^+:Z^+\to B$, its pullback
$(c^+)^*\mathcal A_B$ is the line bundle of the transformed
$mL$. This proves separately that the transformed $L$ is nef
and that $mL$ remains Cartier with the same $m$.

The driving boundary remains klt throughout its LMMP.
The transformed \(\Delta-\tfrac12P\) remains effective and
is smaller than the driving boundary, so it too is klt.
The same estimate and the same integer \(m\) apply inductively
at every subsequent step.
\end{proof}

\subsection{The good-model implication}

We now assume smooth nonvanishing in the current dimension and prove the
good-model implication. The proof first reduces real nef boundaries to
rational ones. The case of positive
Kodaira dimension is then handled by lower-dimensional good models, leaving the
zero-dimensional Iitaka case and its klt perturbations.

\begin{proposition}[Inductive reduction to smooth nonvanishing]
\label{prop:inductive-reduction}
Assume Assumption~\ref{mmp:lower-models}. Assume in addition that
every smooth projective \(n\)-fold with pseudo-effective
canonical divisor has nonnegative Kodaira dimension.
Then every projective lc \(n\)-fold with real boundary and
pseudo-effective adjoint has a good log minimal model.
In particular, every nef rational lc adjoint in dimension \(n\)
is semiample.
\end{proposition}

\begin{proof}
Hashizume's reduction gives nonvanishing and log minimal models
for projective lc pairs with real boundary in dimensions at most
\(n\) \cite[Theorem~1.4]{Hashizume2018}.
Thus we may pass to a \(\Q\)-factorial dlt log minimal model.
Fix the support of its boundary, impose the rational affine
constraints that its coefficient-one components remain equal
to one, and take a sufficiently small rational polytope around
the given boundary within those constraints. On a fixed log
resolution witnessing dlt, the strict discrepancy inequalities
remain strict in this neighborhood; thus its boundaries remain
dlt. The rational-polytope theorem for nef adjoints
\cite[Remark~3.1 and Proposition~3.2(3)]{Birkar2011II}, applied
to all extremal rays and intersected with this neighborhood,
expresses the given real boundary in a finite rational simplex
of dlt boundaries with nef adjoints.
It is enough to prove semiampleness for these
rational adjoints. Their positive real combinations are
semiample, using the product of the associated contractions.
For a rational adjoint, real semiampleness can also be
rationalized: the finite linear equations expressing its
divisor and principal-divisor coefficients have rational data,
so a real solution with positive coefficients has a rational
solution with the same positivity.

We therefore work with a rational dlt pair \((Z,\Delta)\) and
\(L=K_Z+\Delta\) nef.
Real nonvanishing gives rational nonvanishing in this case:
retain the finitely many divisors and principal divisors in
an effective real representative and solve the resulting
rational linear system with nonnegative rational coefficients.
If \(\kappa(Z,L)\geq 1\), lower-dimensional good models give a
good minimal model by \cite[Lemma~3.5]{FujinoGongyoRings2014}.
Lemma~\ref{mmp:comparison} transfers semiampleness back to \(Z\).
It remains to treat
\[
 \kappa(Z,L)=0.
\]

\paragraph{The non-pseudo-effective perturbation.}
Put \(P=\Delta\) in the klt case and
\(P=\lfloor\Delta\rfloor\) in the other case.
Suppose \(L-\epsilon P\) is not pseudo-effective for every
sufficiently small \(\epsilon>0\).
Choose rational \(\epsilon\) as in
Lemma~\ref{mmp:trivial-perturbation}, and run the klt LMMP
with scaling to a Mori fiber space
\[
 (Z,\Delta-\epsilon P)\dashrightarrow
 (Z',\Delta'-\epsilon P')\xrightarrow{f} T.
\]
Existence and termination in the non-pseudo-effective case are
the established klt results of \cite{BCHM2010}.
The program is crepant for \(L\).
The line-bundle descent in
Lemma~\ref{mmp:trivial-perturbation}, applied also to the final
contraction, gives a nef rational divisor \(A\) on \(T\) with
\[
 K_{Z'}+\Delta'\sim_{\Q}f^*A.
\]
The base has dimension less than \(n\).

If \((Z,\Delta)\) is klt, the transformed original pair
\((Z',\Delta')\) is klt as well. Ambro's descent theorem
\cite[Theorem~0.2]{Ambro2005} gives an effective rational
boundary \(\Delta_T\) such that \((T,\Delta_T)\) is klt and
\[
 A\sim_{\Q}K_T+\Delta_T.
\]
All its hypotheses hold: the total pair is globally klt and
effective, \(f\) is a projective contraction, and its adjoint
is rationally linearly pulled back. In particular, this use
requires no conjecture about semiampleness of a moduli divisor.
Assumption~\ref{mmp:lower-models} and
Lemma~\ref{mmp:comparison} make \(A\) semiample.

Otherwise \(P'\) is effective and relatively ample, since
\(K_{Z'}+\Delta'-\epsilon P'\) is relatively antiample.
Some component of the floor therefore dominates \(T\).
Take a crepant dlt blowup $b:(\widetilde Z,\widetilde\Delta)
\to(Z',\Delta')$, and choose the strict transform $S$ of such
a component. The restriction $f_S=(f\circ b)|_S:S\to T$ is
surjective. Crepancy and adjunction give an lc pair on $S$
whose nef adjoint is rationally linearly equivalent to $f_S^*A$.
It is semiample by the lower-dimensional hypothesis.
Semiampleness of a rational line bundle descends along a proper
surjection: use the equality of sections for its connected-fiber
Stein factor, and norms for the remaining finite morphism.
For the latter assertion, choose a section of a basepoint-free
multiple nonzero simultaneously at all points of the chosen
finite fiber. Such a section exists because each vanishing
condition is a proper linear subspace, and finitely many proper
linear subspaces do not cover a complex vector space. Its norm
is nonzero at the point downstairs. Hence \(A\) is semiample
also in this case. Crepancy transfers the conclusion back to \(L\).

\paragraph{The klt pseudo-effective perturbation.}
We next settle the klt case when \(L-\epsilon\Delta\) is
pseudo-effective for some small rational \(\epsilon>0\).
Nonvanishing gives
\[
 L\sim_{\Q}G_0:=H_0+\epsilon\Delta,\qquad H_0\geq 0.
\]
Take a resolution \(p\colon W\to Z\) with reduced SNC divisor
\(D\) consisting of all exceptionals and the strict support of
\(H_0+\Delta\). We have
\[
 K_W+D=p^*L+R,\qquad R\geq 0.
\]
Here every component of \(D\) has positive coefficient in
\[
 G_W:=p^*G_0+R\sim_{\Q}K_W+D.
\]
Indeed, at strict boundary components the difference between
coefficient one and a klt boundary coefficient is positive,
and at exceptional components the coefficient contributed by
the adjoint comparison is the positive log discrepancy.
The pushforward \(p_*G_W\) is bounded coefficientwise by a
fixed multiple of \(G_0\). Section spaces inject under
birational pushforward, so
\[
 0\leq\kappa(W,K_W+D)\leq\kappa(Z,G_0)=0.
\]

Take a \(\Q\)-factorial dlt log minimal model
\((V,D_V)\) of \((W,D)\), and fix a common resolution
\[
 W\xleftarrow{q}U\xrightarrow{v}V.
\]
Its boundary is reduced. On $U$, the comparison
divisor in Lemma~\ref{mmp:comparison} is exceptional over \(V\).
Pushing $q^*G_W$ forward by $v$ therefore gives
\[
 K_V+D_V\sim_{\Q}G_V=\sum_i b_iD_i,\qquad b_i>0,
 \qquad \Supp G_V=D_V.
\]
For completeness, positivity at a component extracted on \(V\)
also holds: its log discrepancy over \((W,D)\) is zero, its
center lies in \(D\), and the pullback of the full-support
effective \(G_W\) has positive order there.
There is no comparison-divisor coefficient at a prime retained
on \(V\).
Effectivity and exceptionality in the comparison preserve the
adjoint section ring, so the nef adjoint on \(V\) still has
Iitaka dimension zero.

The small-perturbation hypothesis of
Theorem~\ref{thm:signed} is now explicit.
If \(D_V\neq 0\), choose rational \(c\) with
\(0<c<\min_i b_i\). Then
\[
 K_V+(1-c)D_V
 \sim_{\Q}\sum_i(b_i-c)D_i\geq 0.
\]
If \(D_V=0\), pseudo-effectivity is immediate.
Proposition~\ref{prop:boundary-restriction} supplies the
required semiampleness on \(D_V\).
Theorem~\ref{thm:signed} thus makes \(K_V+D_V\) abundant.
Its Iitaka dimension is zero, so its numerical dimension is
zero as well. For an ample divisor $H_V$ on $V$, we obtain
\[
 G_V\cdot H_V^{n-1}=(K_V+D_V)\cdot H_V^{n-1}=0.
\]
Effectivity forces $G_V=0$, and its full support then gives $D_V=0$.

Recall that $G_W=p^*G_0+R$ with $R\ge0$. On the same common
resolution,
\[
 0\le q^*p^*G_0\le q^*G_W,
 \qquad v_*q^*G_W=G_V=0.
\]
Thus $q^*p^*G_0$ is effective and $v$-exceptional. It is nef
because it is rationally linearly equivalent to $(p\circ q)^*L$.
The negativity lemma for $v$ forces it to vanish. Therefore
\(G_0=0\) and \(L\sim_{\Q}0\).
Together with the preceding cases, this proves the klt
good-model assertion in dimension \(n\): for a non-nef
pseudo-effective klt adjoint, first take its klt log minimal
model and apply what was just proved.

\paragraph{The remaining non-klt case.}
Finally suppose the original pair is not klt and a small
floor perturbation $L-\delta P$ is pseudo-effective. Choose
rational $\epsilon>0$ with $2\epsilon\le\delta$ and small
enough to keep the perturbed boundaries effective and klt.
Convexity of the pseudo-effective cone, applied between
$L$ and $L-\delta P$, then makes both
\[
 L-\epsilon P,\qquad L-2\epsilon P,
 \qquad P=\lfloor\Delta\rfloor,
\]
pseudo-effective klt adjoints.
Their Iitaka dimensions are zero: nonvanishing gives the lower
bound, and adding the effective perturbation bounds each by
\(\kappa(Z,L)=0\).
The preceding klt case gives good models for both of these
perturbed adjoints.
Consequently their Nakayama numerical dimensions are zero.
Since
\[
 2(L-\epsilon P)
   =L+(L-2\epsilon P)
\]
and the last summand has an effective rational representative,
monotonicity and homogeneity of \(\kappa_\sigma\) give
\[
 \kappa_\sigma(Z,L)
 \leq\kappa_\sigma\bigl(Z,2(L-\epsilon P)\bigr)=0.
\]
The nef divisor \(L\) is therefore numerically trivial.
Nonvanishing is already available in this finishing step:
write \(L\sim_{\Q}E\geq 0\).
For an ample divisor \(H\), the equality
\(E\cdot H^{n-1}=0\) forces \(E=0\).
Thus \(L\sim_{\Q}0\), without any additional nonvanishing or
abundance premise.

This completes the remaining zero-Iitaka-dimension case, so every
rational nef dlt adjoint considered above is semiample. Returning
to the finite rational simplex proves the real-boundary good-model
assertion. Lemma~\ref{mmp:comparison} then descends semiampleness
to every original nef rational lc pair.
\end{proof}
